\documentclass[12pt]{article}

\usepackage{bookmark}
\usepackage{amsmath}
\usepackage[utf8]{inputenc}
\usepackage[T1]{fontenc}
\usepackage{csquotes}
\usepackage{setspace}
\usepackage{graphicx}
\usepackage[a4paper,bottom=2.0cm,top=2.0cm,left=2.0cm,right=2.0cm]{geometry}
\usepackage{float}
\usepackage{titlesec}
\usepackage{color}
\usepackage{caption}
\usepackage{amsthm}
\usepackage{thmtools}
\usepackage{blindtext}
\usepackage{multicol}
\usepackage{lipsum}
\usepackage{wrapfig}
\usepackage{subfig}
\usepackage{hyperref}
\usepackage{cleveref}
\usepackage{listings}

\graphicspath{{Figures},{Plots}}

\newenvironment{Figure}
  {\par\medskip\noindent\minipage{\linewidth}}
  {\endminipage\par\medskip}

\definecolor{bluePoli}{cmyk}{0.4,0.1,0,0.4}
\declaretheoremstyle[
  headfont=\color{bluePoli}\normalfont\bfseries,
  bodyfont=\color{black}\normalfont\itshape,
]{colored}

\captionsetup[figure]{labelfont={color=bluePoli}} % Set colour of the captions
\captionsetup[table]{labelfont={color=bluePoli}} % Set colour of the captions
\captionsetup[algorithm]{labelfont={color=bluePoli}} % Set colour of the captions

\titleformat{\section}
{\color{bluePoli}\normalfont\Large\bfseries}
{\color{bluePoli}\thesection.}{1em}{}
\titlespacing*{\section}
{0pt}{3.3ex}{3.3ex}

\titleformat{\subsection}
{\color{bluePoli}\normalfont\large\bfseries}
{\color{bluePoli}\thesubsection.}{1em}{}
\titlespacing*{\subsection}
{0pt}{3.3ex}{3.3ex}

\titleformat{\subsubsection}
{\color{bluePoli}\normalfont\large\bfseries}
{\color{bluePoli}\thesubsubsection.}{1em}{}
\titlespacing*{\subsubsection}
{0pt}{3.3ex}{3.3ex}

\renewcommand{\title}{Design Task - Lorry Crane Hydraulics}
\renewcommand{\author}{
    Logan Borer\\
    Matteo Sanderson\\
    Arun Thirumani }
\newcommand{\course}{Fluid Power Systems}
\newcommand{\advisor}{Samuel Kärnell}
\newcommand{\firstcoadvisor}{Róger Mateus Sehnem,}
\newcommand{\secondcoadvisor}{Artur Tozzi de Cantuaria Gama}
\newcommand{\YEAR}{2026}

\begin{document}

    \pagenumbering{gobble}
    % DO NOT REMOVE SPACES BETWEEN LINES!

\begin{center}\includegraphics[width=0.6\textwidth]{liu_logo.png}\end{center}

\vspace{-1mm}
\begin{center}\Large{\textbf{\color{bluePoli}{\title}}}\end{center}

\vspace{-0.2cm}
\begin{center}\noindent\fontsize{0.4cm}{0.5cm}\selectfont \bfseries \textsc{\color{bluePoli} Report for Design Task in \course}\end{center}

\vspace{-0.2cm}
\begin{center} \large{\textbf{Authors:}} \end{center}
\begin{center}\noindent\!\!\!\large{\textbf{\author}}\end{center}
\vspace{0.05cm}

\small \normalfont

\vspace{11pt}

\centerline{\rule{1.0\textwidth}{0.4pt}}

\begin{center}
\begin{minipage}[t]{.24\textwidth}
\begin{minipage}{.9\textwidth}
\noindent
\scriptsize{\textbf{Professor:}} \\
\advisor \\
\\
\textbf{Assistant Professor:} \\ 
\firstcoadvisor \\
\secondcoadvisor \\ 
\textbf{Academic Year:} \\
\YEAR \\
\\
\end{minipage}
\end{minipage}% This must go next to `\end{minipage}`
\begin{minipage}{.74\textwidth}
\noindent \textbf{\color{bluePoli} Abstract:} {
    The following report consists of the design and mathematical process which we underwent to design the hydraluic circuit for a lorry crane. The system contains two cylinder, and it was our task to determine the operating conditions for the cylinders and pump, and pick-out the appropriate components for the task. The dimensioning also required us to consider a work-case scenario, which informed our decision during part-selection.
    }
\end{minipage}
\end{center}

\vspace{12pt}
    \setstretch{1.5}

    \newpage
    \tableofcontents
    \newpage

    \pagenumbering{arabic}

    %% Template ends here (barring final line which also needed to be included)

    \section{Design Task Definition}

        The objective of the design is to design a lorry crane, that can operative with a given mass over a defined drive cycle, and withstand a worst-case loading scenario.\\
        For the task a certain values were given, mainly defining the drive cycle. Some values were true for all group, while others were group specific.

        \subsection{General Values}

            The general values for the drive-cycle provided included:\\
            \begin{itemize}
                \item $v_p01$ $\rightarrow$ Normalised velocity vector for piston 1
                \item $v_p02$ $\rightarrow$ Normalised velocity vector for piston 2
                \item $e_1$ $\rightarrow$ Level-arm for cylinder 1
                \item $e_2$ $\rightarrow$ Level-arm for cylinder 2
                \item $\theta_1$ $\rightarrow$ First arm angle
                \item $\theta_2$ $\rightarrow$ Second arm angle
                \item $L_0$ $\rightarrow$ Base height for lorry crane
                \item $p_2$ $\rightarrow$ Pressure on low-pressure side
                \item $\eta$ $\rightarrow$ Oil dynamic viscosity
                \item $c_v$ $\rightarrow$ Laminar leakage loss coefficient
                \item $k_p$ $\rightarrow$ Coloumb friction coefficient
                \item $k_v$ $\rightarrow$ Viscous friction loss coefficient
                \item $k_{\epsilon}$ $\rightarrow$ Displacement coefficient
            \end{itemize}
            The information on the work-case scenario, provided the following values:
            \begin{center}
                \begin{tabular}{cccc}
                    \hline
                    Parameter & & Cylinder 1 & Cylinder 2\\
                    \hline
                    $\theta_{1,wc} [\text{rad}]$ & & 1,2 & 0,35\\
                    $\theta_{2,wc} [\text{rad}]$ & & 0,0  & 0,0\\
                    $e_{1,wc} [\text{m}]$ & & 0,14 & 0,48\\
                    $e_{2,wc} [\text{m}]$ & & 0.57 & 0,17\\
                    \hline
                \end{tabular}
            \end{center}
         
        \subsection{Group-specific values}
            There were many group-specific values that we were given for the design task, these included:
            \begin{itemize}
                \item $M_{t,wc}$ $\rightarrow$ Maximum total load mass
                \item $q_{r,wc}$ $\rightarrow$ Worst-case flow to drive-cycle-max flow ratio
                \item $M_t$ $\rightarrow$ Total load mass
                \item $L_1$ $\rightarrow$ First arm length
                \item $L_2$ $\rightarrow$ Second arm length
                \item $v_{p0}$ $\rightarrow$ Reference cylinder velocity
                \item $t_0$ $\rightarrow$ Drive-cycle Time
                \item $P_{s,max}$ $\rightarrow$ Maximum system pressure
                \item $\eta_c$ $\rightarrow$ Cylinder Efficiency
            \end{itemize}    
            
        \subsection{Denormalisation and discretization}

            From the provided values we can see that the velocity profile which we were provided has been normalised and we have been provided a reference value.
            Using this reference value we denormalised the normalised velocity profile, to determine the effective velocity profile of the two cylinders.
            \begin{align*}
                v_{p1} = v_{p0} * v_{p01}
                v_{p2} = v_{p1} * v_{p02}
            \end{align*}
            Additionally, we discretized time to a number of discrete times equal to the length of the provided vectors for velocity, angles and arms, as such:
            \begin{lstlisting}[language=Octave]
                time = linspace(0,t_0,length(e_1));
            \end{lstlisting}


    \section{Cylinder Sizing and Selection}
        blablabla 
    
    \section{Circuit Drive-cycle Analysis}
        sdfvddsfsdfsdf
        \begin{equation}
            q = C_qA\sqrt{\frac{2}{\rho}\cdot\Delta p}
        \end{equation}
        \begin{figure}[!ht]
            \centering
            \includegraphics[width = .8\linewidth]{liu_logo.png}\\
        \end{figure}

    \section{Pump Selection}

    \section{System Analysis and Efficiency}

    \section{Cooler Dimensioning}
    
\end{document}
